\providecommand\ExamFontPath{../../../字体/}
\documentclass[fontset=exam]{exam-zh}
\usepackage{longtable,booktabs,array,calc,multirow,fvextra}
\usepackage[normalem]{ulem}
\providecommand{\hl}[1]{#1}
\providecommand{\ul}[1]{\uline{#1}}
\DefineVerbatimEnvironment{verbatim}{Verbatim}{breaklines=true,fontsize=\small}
\providecommand{\tightlist}{\setlength{\itemsep}{0pt}\setlength{\parskip}{0pt}}
\newcommand{\pandocbounded}[1]{#1}
\setlength{\emergencystretch}{3em}
\setlength{\tabcolsep}{2pt}
\examsetup{page/size=a4paper}
\title{厦门大学 2022 级 期中考试 A 卷}
\subject[4em]{概率统计}
\begin{document}
\maketitle
\fancyfoot[C]{第 \thepage 页（共 \pageref*{LastPage} 页）}
\pdfbookmark[1]{一、选择题（在各小题四个备选答案中选出一个正确答案，填在题后的括号中，本大题共6个小题，每小题3分，总计18分）}{section4}

\par\medskip{\heiti

一、选择题（在各小题四个备选答案中选出一个正确答案，填在题后的括号中，本大题共 6 个小题，每小题 3 分，总计 18 分）

\par}\smallskip

\begin{enumerate}
\def\labelenumi{\arabic{enumi}.}
\item
  若事件\(A,B\)相互独立，\(P(B) = 0.5\text{，}\allowbreak P(A - B) = 0.3\)，则\(P(B - A) =\)（ ）.
\end{enumerate}

（A）0.1 （B）0.2 （C）0.3 （D）0.5

\begin{enumerate}
\def\labelenumi{\arabic{enumi}.}
\setcounter{enumi}{1}
\item
  设随机变量\(\left. X \right.\sim\left. N(\mu,4^{2})\text{，}\allowbreak Y \right.\sim N(\mu,5^{2})\)，令\(p = P(X \leq \mu - 4)\text{，}\allowbreak q = P(Y \geq \mu + 5)\)，\\
  则（ ）.
\end{enumerate}

（A）\(p > q\) （B）\(p < q\)

（C）\(p = q\) （D）\(p\),\(q\)的大小由\(\mu\)的取值确定

\begin{enumerate}
\def\labelenumi{\arabic{enumi}.}
\setcounter{enumi}{2}
\item
  设\(f_{1}(x)\)为标准正态分布的概率密度，\(f_{2}(x)\)为\(\lbrack - 1,\, 3\rbrack\)上均匀分布的概率密度，\\
  若\(f(x) = \left\{ \begin{array}{l}
  {af_{1}(x),x \leq 0,} \\
  {bf_{2}(x),x > 0,}
  \end{array} \right.\quad(a > 0,b > 0)\)为概率密度，则\(a,b\)应满足（ ）.
\end{enumerate}

（A）\(2a + 3b = 4\) （B）\(3a + 2b = 4\) （C）\(a + b = 1\) （D）\(a + b = 2\)

\begin{enumerate}
\def\labelenumi{\arabic{enumi}.}
\setcounter{enumi}{3}
\item
  设二维随机变量\((X,Y)\)的联合分布律如下：
\end{enumerate}

\begin{longtable}[]{@{}
  >{\centering\arraybackslash}p{(\linewidth - 4\tabcolsep) * \real{0.3333}}
  >{\centering\arraybackslash}p{(\linewidth - 4\tabcolsep) * \real{0.3333}}
  >{\centering\arraybackslash}p{(\linewidth - 4\tabcolsep) * \real{0.3333}}@{}}
\toprule\noalign{}
\endhead
\bottomrule\noalign{}
\endlastfoot
\emph{X Y} & 0 & 1 \\
0 & a & c \\
1 & b & 0.5 \\
\end{longtable}

已知\(P\left\{ Y = 1 \middle| X = 0 \right\} = \frac{1}{2}\text{，}\allowbreak P\left\{ X = 1 \middle| Y = 0 \right\} = \frac{1}{3}\)，则常数\(a,b,c\)分别为（ ）.

（A）\(a = 0.1,b = 0.2,c = 0.1\) （B）\(a = 0.2,b = 0.2,c = 0.1\)

（C）\(a = 0.2,b = 0.1,c = 0.2\) （D）\(a = 0.1,b = 0.2,c = 0.2\)

\begin{enumerate}
\def\labelenumi{\arabic{enumi}.}
\setcounter{enumi}{4}
\item
  若随机事件\(A,B\)的概率分别为\(P(A) = 0.6\text{，}\allowbreak P(B) = 0.5\)，则\(A\)与\(B\)一定（ ）.
\end{enumerate}

（A）相互对立 （B）相互独立 （C）互不相容 （D）相容

\begin{enumerate}
\def\labelenumi{\arabic{enumi}.}
\setcounter{enumi}{5}
\item
  已知存在两组随机变量\(\left. X \right.\sim U( - 1,1)\)与\(Y = X^{2}\)，则\(X\)与\(Y\)之间的关系为（ ）.
\end{enumerate}

\begin{quote}
（A）独立且不相关 （B）不独立但相关

（C）独立但相关 （D）不独立但不相关
\end{quote}

\pdfbookmark[1]{二、填空题（本大题共6小题，每小题3分，总计18分）}{section21}

\par\medskip{\heiti

二、填空题（本大题共 6 小题，每小题 3 分，总计 18 分）

\par}\smallskip

\begin{enumerate}
\def\labelenumi{\arabic{enumi}.}
\item
  据报道，某年龄段亚健康概率约为 1\%，该年龄段中有 10\%的长期不规律作息者，他们亚健康概率为 5\%，则规律作息者的亚健康概率为\underline{\hspace{2cm}}.
\end{enumerate}

\begin{enumerate}
\def\labelenumi{\arabic{enumi}.}
\setcounter{enumi}{1}
\item
  设\(\left. X \right.\sim\left. B(2,p)\text{，}\allowbreak Y \right.\sim B(3,p)\)，若\(P\left\{ X \geq 1 \right\} = \frac{5}{9}\)，则\(P\left\{ Y \geq 1 \right\} =\)\underline{\hspace{2cm}}.
\item
  若\(X\)和\(Y\)相互独立，它们分别服从参数为\(\lambda_{1},\lambda_{2}\)的泊松分布，则\(Z = X + Y\)服从参数为\underline{\hspace{2cm}}的泊松分布.
\item
  设\(P\{X\ge 0,Y\ge 0\}=\frac{3}{7}\)，\(P\{X\ge 0\}=P\{Y\ge 0\}=\frac{4}{7}\)，则\(P\{\max\{X,Y\}<0\}=\)\underline{\hspace{2cm}}.
\item
  设\(X\)和\(Y\)是随机变量，有\(E(X) = 3\text{，}\allowbreak E(Y) = 1\text{，}\allowbreak D(X) = 4\text{，}\allowbreak D(Y) = 9\)，且\(X\)与\(Y\)的相关系数\(\rho_{XY} = 0.25\)，则\(D(5X - 2Y + 10) =\)\underline{\hspace{2cm}}.
\item
  在供暖的季节，住房的平均温度为 20 度，标准差为 2 度，则住房温度与平均温度的偏差的绝对值小于 4 度的概率的下界估值为\underline{\hspace{2cm}}。
\end{enumerate}

\pdfbookmark[1]{三、（10分）一个袋子里有5只球，编号为1，2，3，4，5，现从袋中同时取出3只，用X表示取出的3只球中的最大号码，试求}{section25}

三、（10 分）一个袋子里有 5 只球，编号为 1，2，3，4，5，现从袋中同时取出 3 只，用\(X\)表示取出的 3 只球中的最大号码，试求随机变量\(X\)的分布律和分布函数.

\pdfbookmark[1]{四、（8分）设某机场每天有200架飞机在此降落，任一飞机在某一时刻降落的概率设为0.02，且设各飞机降落是相互独立的. }{section26}

四、（8 分）设某机场每天有 200 架飞机在此降落，任一飞机在某一时刻降落的概率设为 0.02，且设各飞机降落是相互独立的. 试问该机场需配备多少条跑道，才能保证某一时刻飞机需立即降落而没有空闲跑道的概率小于 0.01（每条跑道只能允许一架飞机降落）？

\pdfbookmark[1]{五、（10分）某人乘汽车去火车站乘火车，有两条路可走. 第一条路程较短但交通拥挤，所需时间X（以分钟计）服从N(40，1}{section27}

五、（10 分）某人乘汽车去火车站乘火车，有两条路可走. 第一条路程较短但交通拥挤，所需时间\(X\)（以分钟计）服从\(N(40\text{，}\allowbreak 10^{2})\) ；第二条路程较长，但阻塞少，所需时间\(X\)服从\(N(50\text{，}\allowbreak 4^{2})\). 试问：\\
（1）若动身时离火车开车只有 1 小时，应走哪条路赶上火车的把握大些？\\
（2）又若离火车开车只有 45 分钟，应走哪条路赶上火车把握大些？

\pdfbookmark[1]{六、（12分）设连续型随机变量X的分布函数为：}{section28}

六、（12 分）设连续型随机变量\(X\)的分布函数为：

\[F(x)=\begin{cases}0,&x<0,\\kx^2,&0\le x<1,\\1,&x\ge1.\end{cases}\]，

\begin{quote}
求：（1）常数\(k\)；（2）\(X\)的概率密度；（3）\(X\)落在\(( - 1,\frac{1}{3})\)及\((\frac{1}{2},\frac{3}{2})\)内的概率.
\end{quote}

\pdfbookmark[1]{七、（16分）设平面区域D由曲线y = frac1x及直线y = 0，x = 1，x = e^2所围成，二维随机变量（X}{section31}

七、（16 分）设平面区域\(D\)由曲线\(y = \frac{1}{x}\)及直线\(y = 0\text{，}\allowbreak x = 1\text{，}\allowbreak x = e^{2}\)所围成，二维随机变量\((X,Y)\)在区域\(D\)上服从均匀分布，求：\\
（1）二维随机变量\((X,Y)\)的联合概率密度；\\
（2）\(X\)的边缘概率密度在\(X = 2\)的值.

\pdfbookmark[1]{八、（4分）已知事件X、Y满足XoverlineY = overlineXY，试证明X = Y.}{section32}

八、（4 分）已知事件\(X\text{、}Y\)满足\(X\overline{Y} = \overline{X}Y\)，试证明\(X = Y\).

\pdfbookmark[1]{九、（4分）试证明事件X和Y恰有一个发生的概率为P(X) + P(Y) - 2P(XY).}{section33}

九、（4 分）试证明事件\(X\)和\(Y\)恰有一个发生的概率为\(P(X) + P(Y) - 2P(XY)\).

\end{document}
